\section{Introduction}
\label{sec:intro}

Diffusion models are inspired by nonequilibrium thermodynamics \cite{sohl2015deep} and are now the standard generative
models for pictures \cite{rombach2022latent,betker2023dalle3}, videos \cite{veo3}, 3D protein structures
\cite{abramson2024alphafold3} and many other kinds of data. They reverse a stochastic differential equation
\cite{anderson1982reverse,ho2020denoising,song2021scorebased} using the score
$\mathbf{s}_t(\mathbf{x})=\nabla_{\mathbf{x}}\log p_t(\mathbf{x})$, which can be rewritten as a posterior mean
\cite{robbins1956empirical,vincent2011connection,efron2011tweedie}. Much has been learned about how these models work.
They go through distinct dynamical regimes \cite{biroli2023generative,biroli2024dynamical,achilli2026speciation}, which
become a cascade of speciations when the data have long-range correlations \cite{garnierbrun2026foam}. They generalise
instead of memorising for reasons that are now partly understood
\cite{kadkhodaie2024generalization,achilli2025memorization,bonnaire2025memorize}, and their generalisation can be biased
\cite{garnierbrun2026biased}.

Here we work with Markov sequences, such as time series or one-dimensional physical fields with nearest-neighbour
interactions. We want to understand what a denoiser gains from knowing that the data are Markov, and how to use it. We
start from the Gaussian case, where the score has a linear closed form set by a precision matrix (\cref{sec:gaussian}).
The simplest case where the score is not linear is the Laplace chain, and there we compute the exact score with message
passing \cite{pearl1988probabilistic,mezard2002analytic,mezard2003passing,mezard2009information}
(\cref{sec:tree,sec:grid}). Finally, we learn the kernel (\cref{sec:laplace}) and test the method on the stochastic
Allen--Cahn field \cite{allen1979microscopic} (\cref{sec:allen-cahn}).

Our main contributions are:
\begin{itemize}
\item For any first-order Markov chain, belief propagation gives the exact score, because a Markov chain observed through
  independent noise is a tree (\cref{sec:tree}).
\item When the kernel is unknown, we learn it with expectation--maximisation if we only see noisy sequences, or with
  maximum likelihood if we have the clean ones. The kernel can be a small neural network, and BP turns it into the score.
  In our experiments this is more accurate than denoising score matching with a global MLP or a window CNN, and it needs
  far fewer parameter updates (\cref{sec:laplace}).
\item We discretise the messages on a grid and test the discretisation against the Gaussian closed form. On the Laplace
  chain the grid error is four orders of magnitude below the error of a Gaussian approximation (\cref{sec:grid}).
\item We go beyond the Gaussian and the Laplace chains to a physical field, the stochastic Allen--Cahn equation, where the
  shape of the transition law changes with the state (\cref{sec:allen-cahn}).
\end{itemize}
